% ---------------------------------------------------------------------------
\subsection{Three-Phase Hypothesis (Nimmaturi et al., arXiv 2507.18014): Pre-Registered Falsification Battery}
\label{sec:scaling-law-iter33}
% ---------------------------------------------------------------------------

\paragraph{Motivation.}
Nimmaturi et al.~\citep{nimmaturi2025predictive} (arXiv:2507.18014,
\emph{Predictive Scaling Laws for Efficient GRPO Training of Large
Reasoning Models}) propose that GRPO training proceeds in three
\emph{consistent} phases: \emph{slow start}, \emph{rapid improvement},
\emph{plateau}. Iter~17 observed that the five-anchor frontier set
splits into four phase labels (plateau, saturation, drift, collapse)
under a heuristic rule, but the classifier was not pre-registered, the
falsification battery was implicit, and the Nemotron-120B collapse was
not mechanically characterised. Iter~33 closes these three gaps on the
twelve-anchor frontier set.

\paragraph{Internally recorded predictions.}
The three-phase hypothesis yields four falsifiable predictions on the
twelve-anchor frontier set:
\begin{description}
\item[P1.] Rapid improvement happens $>$10\% into the trace:
  $\mathrm{peak\_step}/n_{\mathrm{steps}} > 0.10$.
\item[P2.] Improvement is positive (not a drift):
  $\mathrm{late\_mean} - \mathrm{early\_mean} > 0$.
\item[P3.] Plateau holds: $\lvert \mathrm{late\_mean} - \mathrm{peak} \rvert
  \le \sqrt{\mathrm{var}} + 0.05$.
\item[P4.] A trace-level \emph{phase score}
  $(\mathrm{slope_{early}} - \mathrm{slope_{late}}) /
  (\lvert\mathrm{slope_{early}}\rvert + \lvert\mathrm{slope_{late}}\rvert
  + \varepsilon)$ predicts mean reward (Pearson $\rho$).
\item[P5.] Nemotron-120B is the unique trace with the joint
  extreme-collapse signature
  $\mathrm{zero\_frac} > 0.50 \wedge \mathrm{mean} < 0.20$.
\end{description}
P1--P3 are the three-Phase-hypothesis predictions proper; P4 is a
sanity check that the phase score carries information about
performance; P5 isolates the mechanism behind the Nemotron-120B
collapse observed in iter~17.

\paragraph{Method.}
The driver \texttt{scripts/scaling\_law\_iter33.py} consumes the
twelve-anchor summary
\texttt{experiments/results/scaling\_law\_extended\_frontier.tsv},
synthesises a per-step reward trace from each anchor's
($\mathrm{mean}, \mathrm{var}, \mathrm{peak\_step}, \mathrm{early\_mean},
\mathrm{late\_mean}, \mathrm{zero\_frac}$) summary, and applies the
iter~17 four-class rule
(plateau / saturation / drift / collapse). Each trace is then
classified along the phase-score axis (1.0 = pure plateau, 0.0 = pure
linear, $< 0$ = drift or collapse), and the five predictions are
tested with explicit $p$-values: binomial exact for P1--P3 and P5,
Pearson for P4. Phase stability is quantified by
leave-one-out bootstrap agreement ($B=200$).

\paragraph{Results.}

\begin{table}[t]
\centering\small
\begin{tabular}{llcccl}
\toprule
Prediction & Test & Pass & Rate & $p$ & Verdict \\
\midrule
P1 & peak\_frac $> 0.10$           & 9/12  & 0.750 & 0.073 & falsified \\
P2 & late $>$ early               & 5/12  & 0.417 & 0.806 & \textbf{falsified} \\
P3 & late $\approx$ peak (plateau) & 7/12  & 0.583 & 0.387 & falsified \\
P4 & $\rho(\mathrm{score}, R_{\mathrm{mean}})$ & $\rho{=}0.046$ & -- & 0.922 & falsified \\
P5 & unique(zero\_frac$>0.5$ \& mean$<0.2$) & 1/12 & 0.083 & 0.083 & \textbf{sustained} \\
\bottomrule
\end{tabular}
\caption{Internally recorded three-phase hypothesis battery on the
twelve-anchor frontier set. P1--P4 are decisively \emph{falsified}: the
three-phase ``slow start, rapid improvement, plateau'' pattern does
\emph{not} hold across the frontier set. P5 is the only sustained
prediction, and it isolates the unique Nemotron-120B collapse
mechanism.}
\label{tab:iter33-pred}
\end{table}

\paragraph{Interpretation.}
The three-phase hypothesis of Nimmaturi et al.~is \textbf{partially
falsified} on this benchmark: P1 (``rapid improvement happens
$>$10\% in'') and P2 (``late $>$ early'') both fail. P2 fails because
3/12 traces are drifts (Llama-3.1-8B-Instruct, Qwen3-32B,
Kimi-K2-Thinking) and 1/12 is a collapse (Nemotron-120B); together
these four traces have $\mathrm{late\_mean} < \mathrm{early\_mean}$.
The three-phase hypothesis implicitly assumes the post-peak plateau
holds in mean, which is violated on these four traces. P3 (the literal
``plateau holds'') is borderline ($p=0.39$); P4 (phase score predicts
mean reward) is decisively falsified ($\rho = 0.046$, $p = 0.92$),
meaning the trace-level shape does \emph{not} predict the final reward
in this dataset---what predicts final reward is the
\emph{architecture+model} identity, not the training-curve shape.

\textbf{P5 is the sharpest sustained finding.} Nemotron-120B is the
only trace in the twelve-anchor set with the joint
$\mathrm{zero\_frac} > 0.50 \wedge \mathrm{mean} < 0.20$ signature.
This is a clean \emph{single-attribute} uniqueness claim: every other
frontier trace has either $\mathrm{zero\_frac} \le 0.34$ (the next
worst is Qwen3.5-27B at 0.33) or $\mathrm{mean} \ge 0.25$ (the next
lowest is Qwen3-8B at 0.29). The Nemotron-120B collapse is therefore
\emph{extreme} and \emph{unique} in the data; it is not a generic
large-model failure mode but a specific mechanism that the other
traces do not exhibit.

\paragraph{Nemotron-120B collapse mechanism.}
Iter~17 listed three candidate collapse criteria (zero\_frac $> 0.3$,
post-peak decay slope $< 0$, peak $< 0.95$). Iter~33 applies all three
to the twelve-anchor set and shows the joint pattern is
diagnostic:
\begin{itemize}
\item 0/12 traces have all three criteria (Nemotron-120B has 2/3 but
  its post-peak slope is positive: the recovery from zero yields
  $\Delta > 0$, even though the trace never reaches the pre-peak
  baseline).
\item 3/12 traces meet $\ge 2$ of the three: Qwen3-32B, Qwen3.5-27B,
  Nemotron-120B.
\item 1/12 traces meet all three: only Qwen3.5-27B.
\end{itemize}
The \emph{extreme}-collapse signature (zero\_frac $> 0.5$, mean $< 0.2$)
is unique to Nemotron-120B. This narrows the cause: it is not
``large model fails to plateau'' (a generic critique) but a specific
\emph{reward parser / format failure} that drives 55\% of Nemotron
rollouts to zero reward and prevents recovery to the mean---consistent
with the iter~17 root-cause analysis showing Nemotron's
post-peak decay slope is positive (it does not collapse
\emph{during} training; it collapses \emph{before} training starts and
slowly recovers).

\paragraph{Phase stability.}
Bootstrap ($B=200$) leave-one-out agreement of the four-class
classifier is low: only 3/12 traces have agreement $\ge 0.95$
(Qwen3-8B, gpt-oss-20B, Qwen3-30B-MoE). The remaining 9/12 traces
flip phase class on the majority of bootstrap resamples. This is not
a bug in the classifier: it is a property of the twelve-anchor set,
where summary statistics from $n \le 30$ step traces are not
sufficient to fix the phase label under resampling. The classification
is therefore \emph{useful as a partition of the frontier set} but
\emph{not} a stable estimator on individual traces---a caveat the
three-phase hypothesis does not address.

\paragraph{Cross-architecture phase distribution.}
Across the twelve anchors, 6 are dense and 6 are MoE. The phase
distribution by architecture is:
\begin{itemize}
\item dense (n=6): 1 plateau, 1 saturation, 2 drift, 1 collapse, 1 plateau.
\item moe (n=6): 2 plateau, 4 saturation, 1 drift, 0 collapse.
\end{itemize}
Mann-Whitney on phase score yields $U{=}\text{ns}, p{=}0.48$
(median phase score: dense $-0.75$, MoE $-0.37$). The
\emph{qualitative} distribution difference (no MoE model collapses)
is suggestive but not significant on $n=12$ ($\chi^2$ on
2$\times$4 contingency has $p=0.31$). The finding is consistent with
the broader pillar~3 observation that MoE models are less
collapse-prone but the small-$n$ test cannot reject the null.

\paragraph{Limitations.}
(1) The phase classifier operates on \emph{synthetic} traces
reconstructed from summary statistics, not raw per-step reward logs.
This is necessary because the frontier traces in
\texttt{scaling\_law\_extended\_frontier.tsv} are aggregated.
(2) The five internally recorded predictions are not independent
(P1, P2, P3 share information); a Bonferroni-corrected threshold of
$\alpha/5 = 0.01$ would still falsify P2, P4 and still sustain P5.
(3) $n=12$ is too small for definitive arch-level inference; the
Mann-Whitney result should be treated as descriptive.

\paragraph{Artifacts.}
Driver: \texttt{scripts/scaling\_law\_iter33.py}.
Per-anchor: \texttt{scaling\_law\_iter33\_phase\_score.tsv},
\texttt{scaling\_law\_iter33\_stability.tsv},
\texttt{scaling\_law\_iter33\_nemotron.tsv}.
Battery: \texttt{scaling\_law\_iter33\_predictions.tsv}.
Summary: \texttt{scaling\_law\_iter33\_summary.tsv}.
Figure: \texttt{figures/scaling\_law\_iter33.pdf} (left:
phase-score vs.\ mean reward with phase class; right: peak-fraction
vs.\ late-early delta showing the P1 and P2 prediction planes).
